\documentclass[11pt]{article}
\usepackage{icmlko}
\begin{document}

\title{두 문이 서로를 닫는다 --- 같은 자료를 공유 열에 넣으면 날짜에서, 전용 열에 넣으면 전이에서 걸린다}
\author{월드모델 실험실}
\date{노트 540 · 2026년 8월 4일}
\maketitle

\begin{abstract}
노트 539 에서 시장팝업 \texttt{category} 를 \textbf{공유 \texttt{grp}} 에 채우니
그 도메인이 두 배($+0.1592$)가 되고 판이 $+0.0048$ 올랐는데 \textbf{날짜 통제
판 $-0.0014$}(노트 446 거부권)로 기각됐다. 하락은 전부 이웃이었으므로,
\textbf{시장팝업 전용 열}로 넣으면 이웃이 안 다칠 것이라 예상했다.

\begin{center}
\begin{tabular}{lrr}
\hline
자 & 공유 \texttt{grp} & \textbf{전용 열} \\
\hline
판 & $+0.0048$ & $+0.0027$ \\
\textbf{날짜 통제 판} & $\mathbf{-0.0014}$ ($0$ 밖) & $\mathbf{-0.0006}$ ($0$ 을 문다) \\
무리 안 판 & $-0.0025$ & $+0.0015$ \\
\textbf{KR 만화} & $-0.0017$ ($0$ 을 문다) & $\mathbf{-0.0060}$ ($0$ 밖) \\
\textbf{비게임 앱} & $-0.0035$ ($0$ 을 문다) & $\mathbf{-0.0050}$ ($0$ 밖) \\
\hline
시장팝업 $\rho$ & $+0.1592$ & $+0.0975$ \\
\hline
\end{tabular}
\end{center}

\textbf{②는 통과했고 ③이 반증됐다.} 전용 열은 날짜 통제를 고치는 대신
\textbf{집 밖 전이를 망친다.} \textbf{둘 다 기각.}
\end{abstract}

\begin{figure}[t]
\centering
\includegraphics[width=\textwidth]{figs/f1.pdf}
\caption{\textbf{가} --- 두 형태의 다섯 자. \textbf{나} --- 덜 벌고 더 잃는다.
\textbf{다} --- 기계.}
\end{figure}

\section{사전 등록 셋의 결과}

\paragraph{① 부분 반증.} 시장팝업 이득 $+0.0975$($12/12$)로 기준
$\ge +0.10$ 에 아슬하게 못 미치고, \textbf{공유 열의 $+0.1592$ 보다 크게
작다.} 같은 정보를 넣었는데 덜 번다.

\paragraph{② 통과.} 날짜 통제 판이 $-0.0014$($0$ 밖 음수)에서
$-0.0006$($4/12$, $0$ 을 문다)으로 나아졌다. \textbf{이웃 손해가 공유 열의
뜻이 바뀌어서라는 읽기가 맞았다} --- 전용 열은 그 부분을 실제로 고친다.

\paragraph{③ 반증.} 그런데 집 밖 짝 둘이 \textbf{$0$ 을 물다가 $0$ 밖
음수로 악화}됐다(KR $-0.0017 \to -0.0060$, 앱 $-0.0035 \to -0.0050$).

\section{기계는 이미 원장에 있었다}

\begin{quote}
\textbf{전용 열은 두 번째 도메인 표지로 작동한다.} 그 열이 채워졌다는 것이
곧 ``이 행은 시장팝업이다''이므로, 나무는 그것을 도메인 지시자로 쓴다.
\end{quote}

노트 $419\!\sim\!423$ 이 같은 것을 이미 쟀다 --- \texttt{grp2} 를 두 번째
표지로 넣으니 KR 전이가 $+0.1946 \to +0.0819$ 로 떨어졌고, 설명 셋(칸 수
$\cdot$ 곱 $\cdot$ 분산)이 차례로 죽은 뒤 남은 처분이 \textbf{``도메인 표지는
하나만 쓴다 --- 왜인지는 모른다''}였다.

\begin{quote}
\textbf{같은 값을 다시 치르고 있다.} 그때는 표지를 \emph{더} 넣어서였고 지금은
전용 열을 넣어서인데, 나무에게는 같은 것이다. \textbf{두 해 전 관측이 새
자리에서 다시 확인됐다.}
\end{quote}

이것은 \textbf{덜 벌고 더 잃는} 이유도 설명한다 --- 전용 열의 정보가
도메인 지시자에 흡수되면, 그 안의 \emph{범주 사이 차이}는 덜 쓰인다
($+0.1592 \to +0.0975$).

\section{두 문이 서로를 닫는다}

\begin{center}
\begin{tabular}{ll}
\hline
넣는 자리 & 걸리는 곳 \\
\hline
공유 \texttt{grp} & 날짜 통제 판(이웃 열 도메인) \\
전용 열 & 집 밖 전이(두 번째 도메인 표지) \\
\hline
\end{tabular}
\end{center}

\textbf{같은 자료인데 어디에 넣느냐로 서로 다른 자에 걸린다.} 자료가 나쁜
것이 아니라 \emph{붙이는 방식}마다 값이 다르다.

\section*{조항}

\begin{quote}
\textbf{한 도메인에만 있는 열은 그 도메인의 지시자다.} 도메인 전용 축을 새로
만들 때는 전이를 먼저 재고, 노트 $419\!\sim\!423$ 의 ``표지는 하나만''을
적용받는다고 본다.
\end{quote}

\section*{남은 길 둘}

\begin{itemize}
\item \textbf{노트 446 거부권 재검토} --- 공유 열 형태의 유일한 실격이
날짜 통제 $-0.0014$ 인데, \textbf{그것이 날짜 구멍이 아님을 노트 539 가
보였다}(category $\leftrightarrow$ 연도 $p = 0.168$, 시장팝업은 그 자에
안 들어간다). \emph{별도 사전 등록}이 필요하다 --- 지금 고치면 점수 고르기다.
\item \textbf{제삼의 형태} --- 이웃 손해를 줄이면서 도메인 지시자가 되지 않는
붙임. 예: \texttt{grp} 에 넣되 시장팝업 행이 결측 갈래에서 빠지는 것을
보정하기.
\end{itemize}

\section*{덧 --- 이 갈래의 값}

기각 둘이지만 \textbf{자료 자체는 살아 있다}. 시장팝업 $\rho$ 를 $+0.157$ 에서
$+0.315$ 로 만드는 한 칸이 있고, 그것을 어떻게 붙일지가 남았을 뿐이다.
\textbf{노트 $530 \to 534 \to 537 \to 538 \to 539 \to 540$ 여섯 노트가
``제일 못 맞히는 도메인''을 ``붙이는 방법만 남은 도메인''으로 바꿨다.}

\end{document}
